\subsection{李子群}

设 G 为李群，H 为 G 的子群，同时也是 G 的子流形。则映射 \((x,y)\mapsto xy^{-1}\) 从 \(H\times H\) 到 G 是解析的，因此映射 \((x,y)\mapsto xy^{-1}\) 从 \(H\times H\) 到 H 也是解析的。（Differentiable and Analytic Manifolds, R, 5.8.5）。因此 H 配以由 G 的群结构和流形结构诱导的结构后是李群。

定义 3. 设 G 为李群。如果 G 的子集 H 是 G 的子群和子流形，则称 H 为李子群。

G 的开子群是 G 的李子群。特别地，如果 G 是实或复李群，则其单位元连通分支是 G 的李子群。

命题 5. 设 G 为李群，H 为 G 的李子群。

\begin{enumerate}
\def\labelenumi{(\roman{enumi})}
\item
  H 在 G 中是闭的。
\item
  H 到 G 的典范嵌入是李群态射。
\item
  设 \(\mathbf{L}\) 为李群，\(f\) 为从 \(\mathbf{L}\) 到 \(\mathbf{G}\) 的映射，满足 \(f(\mathbf{L}) \subset \mathbf{H}\)。要使 \(f\) 成为从 \(\mathbf{L}\) 到 \(\mathbf{H}\) 的态射，充要条件是 \(f\) 是从 \(\mathbf{L}\) 到 \(\mathbf{G}\) 的态射。
\end{enumerate}

由 Differentiable and Analytic Manifolds, R, 5.8.3，H 是局部闭的。因此 H 是闭的（General Topology, Chapter III, § 2, Proposition 4）。断言 (ii) 显然。断言 (iii) 由 Differentiable and Analytic Manifolds, R, 5.8.5 得出。

命题 6. 设 G 为李群，H 为 G 的子群。H 是 G 的李子群的充要条件是：存在一点 \(h \in H\) 以及 G 中 h 的一个开邻域 U，使 \(H \cap U\) 是 G 的子流形。

该条件显然是必要的。设它成立。对于所有 \(h' \in \mathbf{H}\)，平移 \(\gamma (h'h^{-1})\) 是流形 \(G\) 的自同构，并将子流形 \(H \cap U\)（在 \(U\) 中）映入子流形 \((h'h^{-1}H) \cap (h'h^{-1}U)\)（在 \(h'h^{-1}U\) 中）。由于 \(h'h^{-1}H = H\)，且 \(h'h^{-1}U\) 是 \(h'\) 在 \(G\) 中的开邻域，可见 \(H\) 的每一点都有一个开邻域 \(V\)，使 \(V \cap H\) 是 \(G\) 的子流形。因此 \(H\) 是 \(G\) 的子流形。

设 \(G\) 为李群，\(H\) 为 \(G\) 的李子群。若 \(L\) 是 \(H\) 的李子群，则 \(L\) 是 \(G\) 的李子群（由 Differentiable and Analytic Manifolds, \(R\)，5.8.6）。设 \(M\) 是 \(G\) 的李子群且 \(M \subset H\)。则 \(M\) 是 \(H\) 的李子群，因为 \(M\) 到 \(H\) 的典范嵌入显然是浸入。

令 k 为 K 的一个非离散闭子域。G 的李 k-子群是 G 的底层李 k-群的李子群。

备注。若在定义 3 中以 ``子流形'' 替换 ``拟子流形''，便得到 G 的李拟子群的定义。（对于有限维 G，李拟子群恰好就是李子群。）设 K 的特征为 0。对李拟子群，命题 5 以同样的证明仍然成立。命题 6 也以同样的证明成立，只需将 ``李子群'' 替换为 ``李拟子群''，并将 ``子流形'' 替换为 ``拟子流形''。

